\subsection{User-level mitigations}
This gives way to what is known as time protection. Time protection is defined
within \cite{ge2019} as "A collection of OS mechanisms which jointly prevent
interference between security domains that would make execution speed in one
domain dependent on the activities of another". An OS, which implements Time
protection must make it such, that all execution within partitioned parts does
not depend on each other making both covert channels and side channels
infeasible. This is quite a strict version of Time Protection, and to close
covert channels it must be, but quite a significant amount of timing channels
can be closed at the user-level if one allows a distinction between different
information types.

One such example is the FaCT DSL \cite{cauligi2019}. FaCT sought to address the
challenges of writing constant-time cryptographic code, which does not depend
on "secret" information. That is, instead of treating all information as
possible leakage, one indicates what information must not be leaked through a
timing channels, and the FaCT type system guarantees that in isolated execution
this will not leak. FaCT works by making the code constant time in regards to
the sensitive information, and any branching is only dependent on non-sensitive
information. The idea being, that any attack can not measure the execution time
of the encryption code to gather information about sensitive information. This
does work to an extent, by guaranteeing constant time at the C program level,
but compiler introduced and variable execution time of instructions based on
inputs can still leak sensitive information. For example, the division operator
is notoriously known to have its execution dependent on the numerator
\cite{schroder, jin2023}. Mitigations against these attacks would still require
help from hardware, and make use of a specific instructions that guarantee a
constant execution time.

Furthermore, information that the operating system holds about a given process
could also be deemed sensitive. For example, in highly secure environments the
runability of a task could be deemed sensitive, and no changes on the user
level side will ever close such a timing channels, as the OS scheduler will
unknowingly always leak this information through its scheduling decisions.
Therefore, while user level implementations do help significantly, the issue
can not be fully closed without both the support of the OS and hardware.

As such, even though user level mitigations can close some timing channels it
is far from everything and a reset of shared resources through instructions
such as the temporal fence fence\cite{wistoff2023} are still needed even if the
C code itself seems "safe".


